%% LyX 2.2.1 created this file.  For more info, see http://www.lyx.org/.
%% Do not edit unless you really know what you are doing.
\documentclass[english,12pt]{article}
\usepackage[T1]{fontenc}
\usepackage[latin9]{inputenc}
\usepackage{refstyle}
\usepackage{units}
\usepackage{mathtools}
\usepackage{amsmath}
\usepackage{amsthm}
\usepackage{amssymb}
\usepackage{graphicx}

\makeatletter

%%%%%%%%%%%%%%%%%%%%%%%%%%%%%% LyX specific LaTeX commands.

\AtBeginDocument{\providecommand\lemref[1]{\ref{lem:#1}}}
\RS@ifundefined{subsecref}
  {\newref{subsec}{name = \RSsectxt}}
  {}
\RS@ifundefined{thmref}
  {\def\RSthmtxt{theorem~}\newref{thm}{name = \RSthmtxt}}
  {}
\RS@ifundefined{lemref}
  {\def\RSlemtxt{lemma~}\newref{lem}{name = \RSlemtxt}}
  {}


%%%%%%%%%%%%%%%%%%%%%%%%%%%%%% Textclass specific LaTeX commands.
\theoremstyle{plain}
\newtheorem{thm}{\protect\theoremname}[section]
\theoremstyle{definition}
\newtheorem{example}[thm]{\protect\examplename}
\theoremstyle{definition}
\newtheorem{defn}[thm]{\protect\definitionname}
\theoremstyle{plain}
\newtheorem{lem}[thm]{\protect\lemmaname}
\ifx\proof\undefined
\newenvironment{proof}[1][\protect\proofname]{\par
\normalfont\topsep6\p@\@plus6\p@\relax
\trivlist
\itemindent\parindent
\item[\hskip\labelsep\scshape #1]\ignorespaces
}{%
\endtrivlist\@endpefalse
}
\providecommand{\proofname}{Proof}
\fi
\theoremstyle{plain}
\newtheorem{prop}[thm]{\protect\propositionname}
\theoremstyle{plain}
\newtheorem{cor}[thm]{\protect\corollaryname}
\theoremstyle{remark}
\newtheorem{rem}[thm]{\protect\remarkname}

%%%%%%%%%%%%%%%%%%%%%%%%%%%%%% User specified LaTeX commands.
\usepackage[margin=1in]{geometry}

\makeatother

\usepackage{babel}
\providecommand{\corollaryname}{Corollary}
\providecommand{\definitionname}{Definition}
\providecommand{\examplename}{Example}
\providecommand{\lemmaname}{Lemma}
\providecommand{\propositionname}{Proposition}
\providecommand{\remarkname}{Remark}
\providecommand{\theoremname}{Theorem}

\begin{document}

\title{Math 525: Lecture 20}

\date{March 27, 2018}
\maketitle

\section{Limiting distribution}

We have already seen many situations involving the \emph{limiting
distribution}
\begin{equation}
\alpha^{\intercal}=\lim_{n\rightarrow\infty}\mu^{\intercal}P^{n}\label{eq:limiting_distribution}
\end{equation}
of a Markov chain, where $\mu$ is the \emph{initial distribution}
(i.e., $\mathbb{P}(X_{0}=i)=\mu_{i}$). To stress that $\alpha$ is
a function of $\mu$, we will write $\alpha(\mu)$.
\begin{example}[Gambler's ruin]
We determined the probability of ruin by substituting
\[
P=\begin{pmatrix}1 & 0\\
\nicefrac{1}{2} & 0 & \nicefrac{1}{2}\\
 & \nicefrac{1}{2} & 0 & \nicefrac{1}{2}\\
 &  & \ddots & \ddots & \ddots\\
 &  &  & \nicefrac{1}{2} & 0 & \nicefrac{1}{2}\\
 &  &  &  & 0 & 1
\end{pmatrix}\qquad\text{and}\qquad\mu=\begin{pmatrix}0\\
\vdots\\
0\\
1\\
0\\
\vdots\\
0
\end{pmatrix}
\]
into (\ref{eq:limiting_distribution}). There, we found that
\[
\alpha(\mu)=\begin{pmatrix}\nicefrac{1}{2}\\
0\\
\vdots\\
0\\
\nicefrac{1}{2}
\end{pmatrix}.
\]
\end{example}
More generally, it is natural to ask:
\begin{enumerate}
\item Is (\ref{eq:limiting_distribution}) always independent of the initial
distribution $\mu$?
\item Is (\ref{eq:limiting_distribution}) an equilibrium/stationary distribution?
(i.e., $\alpha(\mu)^{\intercal}P=\alpha(\mu)^{\intercal}$)
\end{enumerate}
The answer to the first question is very easily seen to be no:
\begin{example}
Consider the transition matrix $P=I$. Then, $P^{n}=I$ and hence
\[
\alpha(\mu)=\lim_{n\rightarrow\infty}\mu^{\intercal}P^{n}=\lim_{n\rightarrow\infty}\mu^{\intercal}=\mu^{\intercal}.
\]
\end{example}
Worse yet, the limit does not even have to exist:
\begin{example}
Consider
\[
P=\begin{pmatrix}0 & 1\\
1 & 0
\end{pmatrix}\qquad\text{and}\qquad\mu=\begin{pmatrix}1\\
0
\end{pmatrix}.
\]
Note that $P^{2n}=I$ and $P^{2n+1}=P$. Therefore,
\[
\mu^{\intercal}P^{2n+1}=\begin{pmatrix}0 & 1\end{pmatrix}\qquad\text{and}\qquad\mu^{\intercal}P^{2n}=\begin{pmatrix}1 & 0\end{pmatrix}
\]
so that $\alpha(\mu)$ does not exist.
\end{example}
As for the second question, suppose the limit exists. Then, by continuity,
\[
\alpha(\mu)^{\intercal}P=\left(\lim_{n\rightarrow\infty}\mu^{\intercal}P^{n}\right)P=\lim_{n\rightarrow\infty}\mu^{\intercal}P^{n+1}=\alpha(\mu)^{\intercal}.
\]

For the remainder, we will try to understand when the limit is independent
of the initial distribution. As in the gambler's ruin, the limiting
distribution tells us a lot about the problem, and can even be considered
as an approximation for $\mu^{\intercal}P^{n}$ with $n$ large. Moreover,
we will see that in many of these cases, the limit can be computed
efficiently.

\section{Primitive matrices}
\begin{defn}
Let $A$ be a square matrix that is nonnegative (i.e., $A_{ij}\geq0$).
We say $A$ is \emph{primitive }if there exists a positive integer
$m$ such that $A^{m}$ is positive (i.e., $(A^{m})_{ij}>0$).
\end{defn}
\begin{lem}
\label{lem:primitive_implies_reducible}If $A$ is primitive, then
$A$ is irreducible.
\end{lem}
\begin{proof}
We prove this by contrapositive. Suppose $A$ is reducible. By definition,
we can find a permutation matrix $K$ such that
\[
KAK^{\intercal}=\begin{pmatrix}B & C\\
0 & D
\end{pmatrix}.
\]
Let $m$ be a positive integer. Then,
\[
\left(KAK^{\intercal}\right)^{m}=KA^{m}K^{\intercal}=K\begin{pmatrix}B^{m} & X\\
0 & D^{m}
\end{pmatrix}K^{\intercal}
\]
where $X$ is some matrix (depending on $B$, $C$, $D$, and $m$).
The above implies that $A$ is not primitive since $A^{m}$ has some
entries equal to zero.
\end{proof}
\begin{prop}
\label{prop:primitive_characterization}Let $A$ be a nonnegative
square matrix. Then, $A$ is irreducible and aperiodic (i.e., has
period equal to one) if and only if it is primitive.
\end{prop}
Before we give a proof, let's discuss some consequences. Given a matrix
$A$, recall that
\[
i\rightarrow j
\]
means that we can find a walk
\[
i=i_{1}\dashrightarrow i_{2}\dashrightarrow\cdots\dashrightarrow i_{k}=j
\]
from $i$ to $j$. We refer to the number of edges in a walk as its
\emph{length} (in the above example, the length is $k-1$). Recall
also that the period of $i$ is defined as
\[
d(i)=\operatorname{gcd}(\left\{ n\geq1\colon(A^{n})_{ii}>0\right\} ).
\]
Moreover,
\[
(A^{n})_{ii}>0\iff\text{there exists a walk of length }n\text{ from }i\text{ to itself}.
\]
Therefore, we can equivalently define the period as
\[
d(i)=\operatorname{gcd}(\left\{ \operatorname{length}(w)\colon w\text{ is a walk from }i\text{ to itself}\right\} ).
\]
This observation gives us a quick way to check if a matrix is primitive
by checking the gcd of walks from a node to itself. One particularly
useful consequence is given below.
\begin{cor}
\label{cor:quick}Let $A=(A_{ij})$ be a nonnegative and irreducible
square matrix. If $A_{ii}>0$ for some $i$, then $A$ is primitive.
\end{cor}
\begin{proof}
By the assumptions, $d(i)=1$ and hence $A$ is aperiodic. Therefore,
by Proposition \ref{prop:primitive_characterization}, $A$ is primitive.
\end{proof}
\begin{example}
The matrix
\[
P=\begin{pmatrix}\nicefrac{1}{2} & \nicefrac{1}{2} & 0\\
0 & 0 & 1\\
1 & 0 & 0
\end{pmatrix}.
\]
satisfies Corollary \ref{cor:quick}. Indeed, matrix multiplication
reveals
\[
P^{4}=\begin{pmatrix}0.5625 & 0.3125 & 0.1250\\
0.2500 & 0.2500 & 0.5000\\
0.6250 & 0.1250 & 0.2500
\end{pmatrix}.
\]
\end{example}
Let's now return to our goal of proving Proposition \ref{prop:primitive_characterization}.
First, we need an intermediate result from number theory/algebra,
which we state without proof.
\begin{lem}[Semigroup lemma]
\label{lem:semigroup}Any set of non-negative integers which is closed
under addition and which has greatest common divisor 1 must contain
all but finitely many of the non-negative integers.
\end{lem}
We will just prove the forward direction of Proposition \ref{prop:primitive_characterization}:
\begin{proof}[Proof of Proposition \ref{prop:primitive_characterization} ($\Rightarrow$)]
 Let $A$ be an irreducible and aperiodic matrix. For each row $i$
of $A$ let 
\[
J_{i}=\left\{ n\geq1\colon(A^{n})_{ii}>0\right\} .
\]
Now, let $n_{1}$ and $n_{2}$ be elements of $J_{i}$. Then,
\[
(A^{n_{1}+n_{2}})_{ii}=(A^{n_{1}}A^{n_{2}})_{ii}=\sum_{k}(A^{n_{1}})_{ik}(A^{n_{2}})_{ki}\geq(A^{n_{1}})_{ii}(A^{n_{2}})_{ii}>0
\]
and hence $n_{1}+n_{2}$ is also in $J_{i}$. That is, $J_{i}$ is
closed under addition. Moreover, $\operatorname{gcd}(J_{i})=1$ by
the presumed aperiodicity.

Applying \lemref{semigroup} to $J_{i}$, we can find $M(i)$ such
that for all $n\geq M(i)$, we have $(A^{n})_{ii}>0$. Since $A$
is irreducible, we can find $m(i,j)$ such that $(A^{m(i,j)})_{ij}>0$.
Therefore, for $n\geq M(i)$,
\[
(A^{n+m(i,j)})_{ij}=\sum_{k}(A^{n})_{ik}(A^{m(i,j)})_{kj}\geq(A^{n})_{ii}(A^{m(i,j)})_{ij}>0.
\]
Let $M=\max_{i}\{M(i)\}+\max_{i,j}\{m(i,j)\}$. Then, $(A^{M})_{ij}>0$
for all $i,j$ as desired.
\end{proof}
\begin{defn}
A Markov chain whose transition matrix is primitive is called \emph{regular}.
\end{defn}
\begin{prop}
Let $P$ be the transition matrix of a regular Markov chain. Then,
there exists a vector $\alpha$ such that for any vector $\mu$,
\begin{equation}
\alpha^{\intercal}=\lim_{n\rightarrow\infty}\mu^{\intercal}P^{n}.\label{eq:limit}
\end{equation}
Moreover,
\begin{equation}
\alpha^{\intercal}P=\alpha^{\intercal}.\label{eq:equilibrium}
\end{equation}
In fact, $\alpha=c_{1}v_{1}$ where $v_{1}$ is a positive eigenvector
corresponding to the eigenvalue $\lambda=1$ of $P^{\intercal}$ and
$c_{1}$ is a normalizing constant which ensures that $\alpha$ is
a probability vector.
\end{prop}
You are asked to prove this when $P^{\intercal}$ admits a full set
of linearly independent eigenvectors in assignment 8. The general
case can also be proved using the Jordan decomposition of $P^{\intercal}$.
Some observations:
\begin{itemize}
\item (\ref{eq:limit}) states that the limiting distribution is independent
of $\mu$.
\item (\ref{eq:equilibrium}) states that the limiting distribution is an
equilibrium of the Markov chain.
\item $\alpha$ can be obtained by computing an eigenvector of $P^{\intercal}$
associated with the eigenvalue $\lambda=1$. For small matrices, you
can do this by hand. In practice, there are various computational
methods to do this, the simplest of which is the \emph{power method}.
\end{itemize}
Of course, there are ``irregular'' Markov chains which still have
limiting distributions:
\begin{example}
Note that
\[
\mu^{\intercal}\begin{pmatrix}1 & 0\\
1 & 0
\end{pmatrix}=\begin{pmatrix}\mu_{1}+\mu_{2}\\
0
\end{pmatrix}^{\intercal}
\]
for any vector $\mu$.
\end{example}

\section{Page Rank}

We close by discussing an application of regular Markov chains: Page
Rank. This is an algorithm used by Google to determine the ``goodness''
of a page so as to determine how to order the results of a search
query.

\begin{figure}
\begin{centering}
\includegraphics[width=3in]{page_rank}
\par\end{centering}
\caption{The ``goodness'' of each page is proportional to the number of links
to that page.}
\end{figure}

Consider a web with $N$ pages. Links between pages are recorded in
a binary matrix $G=(G_{ij})$. That is,
\[
G_{ij}=\begin{cases}
1 & \text{if there is a link from page }i\text{ to page }j\\
0 & \text{otherwise.}
\end{cases}
\]
Letting $e$ denote the column vector of ones in $\mathbb{R}^{N}$,
note that
\[
(Ge)_{i}=\text{number of outgoing links from }i.
\]

We wish to model a \emph{web surfer}. Suppose first that each page
has at least one outgoing link. If the surfer is at page $i$ at time
$n$, then they choose a page from the set
\[
\left\{ j\colon G_{ij}=1\right\} 
\]
to visit at time $n+1$ with uniform probability. Our assumption guarantees
that the above is nonempty. The transition matrix corresponding to
this Markov chain is
\[
\tilde{P}=\operatorname{diag}(Ge)^{-1}G
\]
where
\[
\operatorname{diag}(x)=\begin{pmatrix}x_{1}\\
 & x_{2}\\
 &  & \ddots\\
 &  &  & x_{N}
\end{pmatrix}
\]
is the diagonal matrix obtained by placing the entries of the vector
$x$ on the diagonal. There are two problems with this preliminary
model: (1) it's not realistic and (2) it's not a regular Markov chain.

We would like to incorporate the notion that the web surfer is \emph{restless}
and can, at any point in time, pick any page (not necessarily via
an outgoing link) at random. Let
\[
E=ee^{\intercal}=\begin{pmatrix}1 & 1 & 1 & \cdots & 1\\
1 & 1 & 1 & \cdots & 1\\
1 & 1 & \vdots &  & \vdots\\
1 & 1 & 1 & \cdots & 1
\end{pmatrix}
\]
where $e$ is the vector of ones. Then, incorporating the restlessness,
we arrive at the new transition matrix 
\[
P=\left(1-\alpha\right)\operatorname{diag}(Ge)^{-1}G+\alpha\frac{1}{m}E
\]
where $0<\alpha\leq1$. The closer $\alpha$ is to one, the ``more
restless'' our web surfer. Since $\frac{1}{m}E$ is a positive transition
matrix, $P^{\prime}$ is trivially a regular Markov chain.
\begin{rem}
The assumption that each page has at least one outgoing link is not
particularly restrictive. Indeed, we can replace $G$ by $G^{\prime}$,
which is defined by
\[
G_{ij}^{\prime}=\begin{cases}
G_{ij} & \text{if }(Ge)_{i}>0\\
1 & \text{otherwise}.
\end{cases}
\]
This corresponds to a surfer that, having encountered a page with
no outgoing links, becomes restless and picks a page at random.
\end{rem}
Since $P$ is regular, the limiting distribution
\[
\nu^{\intercal}=\lim_{n\rightarrow\infty}\mu^{\intercal}P^{n}
\]
gives us an idea of the random surfers distribution after $n$ steps,
assuming $n$ is large. Since $P$ is regular, this does not depend
on the initial distribution $\mu$. Now, let
\[
\hat{\nu}=\begin{pmatrix}1 & \nu_{1}\\
2 & \nu_{2}\\
\vdots & \vdots\\
N & \nu_{N}
\end{pmatrix}
\]
and sort the rows of $\hat{\nu}$ according to the second column:
\begin{example}
Suppose
\[
\nu=\begin{pmatrix}0.2\\
0.1\\
0.3\\
0.4\\
0.1
\end{pmatrix}\implies\hat{\nu}=\begin{pmatrix}1 & 0.2\\
2 & 0.1\\
3 & 0.3\\
4 & 0.4\\
5 & 0.1
\end{pmatrix}.
\]
After sorting,
\[
\hat{\nu}_{\text{sorted}}=\begin{pmatrix}4 & 0.4\\
3 & 0.3\\
1 & 0.2\\
2 & 0.1\\
5 & 0.1
\end{pmatrix}
\]
and hence we can conclude that page 4 has the highest rank, page 3
the second highest, etc.
\end{example}

\end{document}
